\section{Method: FORGE}
\label{sec:method}

We call the proposed protocol FORGE (\textbf{For}ward-forgetting via
decayed aggregation and pre-committed erasure summaries). It has three
components: (i) a \emph{feature-decoupled decayed-aggregation training loop}
that provides forward forgetting on the shared representation; (ii) a
\emph{pre-committed differentially private summary channel} through which
each client transmits a small digest of their local data at training time;
and (iii) a \emph{non-interactive erasure procedure} that, at request time,
uses the departing client's summary to subtract the memorised residual from
the classifier head. The split between feature and head parameters is
essential for the composition theorem (Theorem~\ref{thm:composition}): decay
operates on features, subtraction operates on the head, and the two
mechanisms therefore have independent Gaussian noise sources.

Partition the model's parameters as $\theta = (\theta_\mathrm{sh},
\theta_\mathrm{cl})$ where $\theta_\mathrm{sh}$ are the \emph{shared
representation} parameters (all layers except the final classifier) and
$\theta_\mathrm{cl}$ are the \emph{classifier head} parameters (the final
linear layer that maps features to class logits).

\subsection{Component 1: Feature-decoupled decayed aggregation}
\label{sec:method-decay}

At each round $t$, each sampled client $i \in S^{(t)}$ returns an update
$\Delta_i^{(t)} = (\Delta_{i,\mathrm{sh}}^{(t)}, \Delta_{i,\mathrm{cl}}^{(t)})$
clipped so that the full vector has L2 norm at most $C$. The server
maintains a running exponential moving average (EMA) aggregation state
$\tilde\Delta^{(t)}$ on the feature block only:
\begin{equation}
\tilde\Delta_\mathrm{sh}^{(t+1)}
  = \alpha\, \tilde\Delta_\mathrm{sh}^{(t)}
  + (1-\alpha) \cdot
     \tfrac{1}{|S^{(t)}|} \sum_{i \in S^{(t)}} \Delta_{i,\mathrm{sh}}^{(t)}.
\label{eq:decay-state}
\end{equation}
The feature block is updated with calibrated noise:
\begin{equation}
\theta_\mathrm{sh}^{(t+1)}
  = \theta_\mathrm{sh}^{(t)}
  + \eta_g \cdot \big(\tilde\Delta_\mathrm{sh}^{(t+1)} + \xi^{(t)}\big),
\qquad \xi^{(t)} \sim \mathcal{N}(0, \sigma^2 I_{d_\mathrm{sh}}).
\label{eq:feature-update}
\end{equation}
The head block is updated with unweighted FedAvg and no server noise:
\begin{equation}
\theta_\mathrm{cl}^{(t+1)}
  = \theta_\mathrm{cl}^{(t)}
  + \eta_g \cdot
     \tfrac{1}{|S^{(t)}|} \sum_{i \in S^{(t)}} \Delta_{i,\mathrm{cl}}^{(t)}.
\label{eq:head-update}
\end{equation}
This asymmetry is the architectural anchor of the composition theorem.
The feature block is protected during training by the combined decay and
Gaussian noise mechanism; the head block is updated without noise and the
memorised residual is removed at request time by the summary-subtraction
step.

\subsection{Component 2: Pre-committed summary channel}
\label{sec:method-summary}

Alongside their gradient update, each client transmits, every $F$ rounds, a
differentially private summary $\sigma_i$ of their local data. The default
summary is class-conditional feature moments:
\begin{align}
\sigma_i^{\mathrm{mom}} &= \{(\mu_i^c + \eta^c)\}_{c=1}^{C_\mathrm{class}},\quad \eta^c \sim \mathcal{N}(0, \sigma_\summary^2 I), \nonumber \\
\mu_i^c &= \tfrac{1}{n_i^c} \sum_{j: y_j = c} \phi(x_j).
\label{eq:summary-moments}
\end{align}
Here $\phi(x_j)$ is the clipped feature embedding under the current global
model, and the per-class means are clipped to L2 norm $\le 1$ before
privatising. The pre-noise clipping fixes the summary sensitivity at
$\Delta = 1$ for every release independent of the feature-block state
$\theta_\mathrm{sh}$, which is what makes the $\mu_\summary = \Delta /
\sigma_\summary$ parameter in Theorem~\ref{thm:composition} constant across
training. The Gaussian noise scale $\sigma_\summary$ is calibrated to
give $(\eps_\summary, \delta_\summary)$-DP per summary release against the
client's data; the ablation (Section~\ref{sec:results-ablation}) sweeps
$\sigma_\summary$ at fixed $\Delta$.

Three alternative summary types are implemented in the released code: distilled logits
on a public probe set, DP synthetic microbatches \cite{harder2021dpmerf}, and
random-projection sketched gradients. The ablation across these variants is identified
as future work in Section~\ref{sec:limitations}.

Each summary is committed to a server-side hash-indexed store at the time it is
received. The server retains two artefacts per client, namely (i) the hash of every
released summary (for the verification transcript described below) and (ii) the single
most-recent full summary $\sigma_i$. On receipt of a new release, the previous full
summary is purged. This storage model guarantees that, at erasure time, only one
differentially private release of $\sigma_\forget$ is observable by any party, and the
summary-channel DP cost does not compose across releases. The server never uses
summaries during training; they are accessed only at erasure time.

\subsection{Component 3: Non-interactive erasure}
\label{sec:method-erasure}

When an erasure request arrives for client $i^\star$ at round
$r_\mathrm{req}$, the server performs Algorithm~\ref{alg:erasure}:

\begin{enumerate}[leftmargin=*,itemsep=1pt]
\item Retrieve $\sigma_{i^\star}$ and verify its hash against the client's commitment.
\item Compute the residual direction $v_{i^\star} = \nabla_{\theta_\mathrm{cl}} \mathcal{L}_\mathrm{sub}(\theta_\mathrm{cl};\, \sigma_{i^\star})$, where $\mathcal{L}_\mathrm{sub}$ treats each per-class summary mean as a labelled virtual example under the head's cross-entropy loss.
\item Apply a bounded ascent step on the head only: $\theta_\mathrm{cl}' = \theta_\mathrm{cl} + \gamma\, \Pi_{B(0, R)}(v_{i^\star})$, where $\gamma$ is a step size and $\Pi_{B(0, R)}$ projects onto an $L_2$-ball of radius $R$.
\item Compute $T = r_\mathrm{req} - r_\forget$ and emit a certificate $\mathrm{cert} = (\mu_1(T), \mu_\summary, \mu_\mathrm{joint}(T), \delta_\summary, \mathrm{hash}(\sigma_{i^\star}), \mathrm{tr})$, with $\mu_1(T)$ from Theorem~\ref{thm:forward-forgetting}, $\mu_\summary$ from Lemma~\ref{lem:subtraction}, and $\mu_\mathrm{joint}(T) = \sqrt{\mu_1(T)^2 + \mu_\summary^2}$ from Theorem~\ref{thm:composition}.
\item Purge $\sigma_{i^\star}$ from server storage.
\end{enumerate}

The feature block is not modified at erasure time. By
Theorem~\ref{thm:forward-forgetting}, the feature block's residual dependence on
$i^\star$'s data is already bounded at level $\mu_1(T)$ by the decayed-aggregation
training itself; the head subtraction cleans the remaining dependence localised in the
classifier.

\paragraph{Verification protocol.}
The certificate emitted at step~4 of Algorithm~\ref{alg:erasure} is
$\mathrm{cert} = (\eps_\mathrm{total}, \delta_\mathrm{total}, h_\sigma, \mathrm{tr})$.

\textit{Hash and commitment primitives.} $h_\sigma = H(\forget \| r_s \|
\sigma_\forget)$ where $H$ is a cryptographic hash function (we instantiate with
SHA-256). The client's public commitment, published at summary release time, is a
Pedersen commitment $\mathrm{Com}(h_\sigma, r) = g^{h_\sigma} h^{r}$ in a
prime-order cyclic group $\mathbb{G}$ of order $q$ (we instantiate with
ristretto255), with generators $g, h \in \mathbb{G}$ for which $\log_g h$ is
unknown, and $r$ a uniform blinding scalar in $\mathbb{Z}_q$ sampled by the client.
The client retains $r$ locally; losing $r$ means the client cannot open the
commitment and therefore cannot be served by the erasure procedure, so $r$ must
be persisted alongside the client's own data for the lifetime of the deployment.
Computational binding of the commitment holds under the discrete-log assumption on
$\mathbb{G}$; perfect hiding holds unconditionally. Collision resistance of $H$
holds in the standard model against probabilistic polynomial-time adversaries,
which prevents the server from substituting a different $\sigma_\forget$ with the
same $h_\sigma$ after the fact.

\textit{Protocol parameters.} $\gamma$ and $R$ are public protocol parameters
fixed before training begins (part of the published forward-forgetting schedule);
they are not server-selectable per erasure. Any attempt by the server to use
different values is detectable because the authority recomputes the head delta
using the published $\gamma, R$ and compares the resulting hash against
$h_\mathrm{post}$ in the transcript.

\textit{Transcript.} $\mathrm{tr} = (r_\mathrm{req}, T, h_\mathrm{pre},
h_\mathrm{post})$ with $h_\mathrm{pre} = H(\theta_\mathrm{cl}^{(r_\mathrm{req})})$
and $h_\mathrm{post} = H(\theta_\mathrm{cl}^{(r_\mathrm{req})\prime})$.

\textit{Trust model for the authority.} The supervisory authority is modelled as
semi-honest: it follows the verification protocol and does not share information
with any party. The authority is granted read access to
$\theta_\mathrm{cl}^{(r_\mathrm{req})}$ (the pre-erasure head weights) under a
non-disclosure obligation, so it can recompute the ascent direction $v$ from the
opened summary $\sigma_\forget$ and verify
$H(\theta_\mathrm{cl}^{(r_\mathrm{req})} + \gamma \Pi_{B(0,R)}(v)) = h_\mathrm{post}$.
Under this trust model, the Chen-et-al unlearning-leakage channel
\cite{chen2021unlearningleaks} is closed with respect to \emph{third parties}
(who see only $h_\mathrm{pre}, h_\mathrm{post}$ and the certificate) but not with
respect to the authority itself. For deployments that require a stronger model
(e.g., the authority is itself distrusted), the verification can be replaced by a
succinct zero-knowledge proof of the hash-chain relation on
$(\theta_\mathrm{cl}, v, \gamma, R)$; this substitution is orthogonal to the
privacy analysis and is deferred to future work.

\textit{Authority verification.} The authority (a) opens the Pedersen commitment
with $(h_\sigma, r)$ provided by the client and checks it against the stored
commitment; (b) recomputes $\mu_\mathrm{joint}(T) = \sqrt{\mu_1(T)^2 + \mu_\summary^2}$
from the published forward-forgetting schedule using the certificate's $T$ and
derives the corresponding $(\eps_\mathrm{total}, \delta_\mathrm{total})$ via the GDP
duality; (c) recomputes $v$ from $\sigma_\forget$ and the pre-erasure head, then
verifies $H(\theta_\mathrm{cl}^{(r_\mathrm{req})} + \gamma \Pi_{B(0,R)}(v)) =
h_\mathrm{post}$.

\paragraph{Composition over sequential erasures.}
Consider $K$ clients $\forget_1, \ldots, \forget_K$ that depart in sequence with
cooling-off periods $T_1, \ldots, T_K$, each triggering a head subtraction using
their own pre-committed summary $\sigma_{\forget_k}$. The ascent step for client
$\forget_k$ depends on $\sigma_{\forget_k}$ only, and summaries of different
clients are independent by construction (each is a Gaussian release on the
submitting client's data alone). Two distinct guarantees follow.

\textit{Per-client guarantee.} Each client's data is observed through exactly one
Gaussian mechanism (its own summary release), independent of how many other
clients subsequently request erasure. The per-client $(\eps_\summary,
\delta_\summary)$-DP bound (equivalently $\mu$-GDP with $\mu = \Delta/\sigma_\summary$)
is therefore invariant in $K$: for any individual client $\forget_k$, an adversary
observing the full $K$-erasure history gains no additional information about
$\forget_k$'s data beyond what the single release $\sigma_{\forget_k}$ exposes.

\textit{Joint-dataset guarantee.} For an adversary interested in the joint
distribution of all $K$ clients' data, the $K$ Gaussian releases are independent
and the composed bound follows from Gaussian-DP composition
\cite{dong2019gaussian}: the joint release satisfies $\mu_\mathrm{joint}$-GDP with
$\mu_\mathrm{joint} = \mu \sqrt{K}$. At $\mu = 0.413$ and $K = 100$ handled
erasures, $\mu_\mathrm{joint} \approx 4.1$, which converts to a joint
$(\varepsilon, \delta)$-DP budget of approximately $(4.6, 10^{-5})$; at $K = 10$,
$\mu_\mathrm{joint} \approx 1.3$ and the joint budget is approximately
$(2.5, 10^{-5})$.

The feature-block guarantee retains the forward-forgetting rate for every
departed client whose cooling-off period $T_k \ge T^\star$ at request time, by
the same independence-of-noise argument used in Theorem~\ref{thm:composition}.

The $K$ head-subtraction steps compose on $\theta_\mathrm{cl}$ without inflating any
individual client's budget: the $k$-th subtraction is a deterministic post-processing
of $\sigma_{\forget_k}$ applied at the current head state
$\theta_\mathrm{cl}^{(k-1)}$, so by DP post-processing its cost against
$D_{\forget_k}$ is $\mu_\summary$ regardless of the prior $k{-}1$ modifications to the
head. The chain of head states $\theta_\mathrm{cl}^{(0)}, \ldots,
\theta_\mathrm{cl}^{(K)}$ is a sequence of post-processings of
$(\sigma_{\forget_1}, \ldots, \sigma_{\forget_k})$, which does not alter the per-client
guarantee.

\begin{algorithm}[t]
\footnotesize
\caption{FORGE training round $t$.}
\label{alg:decayed-aggregation}
\begin{algorithmic}[1]
\Require $\alpha$, $C$, $\sigma$, $\eta_g$, $F$, $(\eps_\summary, \delta_\summary)$; state $\tilde\Delta_\mathrm{sh}^{(t)}$, $\theta^{(t)}$
\State Sample $S^{(t)}$; broadcast $\theta^{(t)}$
\ForAll{$i \in S^{(t)}$}
  \State $\Delta_i \gets \mathrm{clip}(\mathrm{LocalTrain}(\theta^{(t)}, D_i^{(t)}), C)$
  \If{$t \bmod F {=} 0$} send $\sigma_i, \mathrm{hash}(\sigma_i)$ \EndIf
\EndFor
\State $\tilde\Delta_\mathrm{sh}^{(t+1)} \!\gets\! \alpha \tilde\Delta_\mathrm{sh}^{(t)} + (1{-}\alpha) \overline{\Delta_\mathrm{sh}}$ \Comment{feature decay}
\State $\theta_\mathrm{sh}^{(t+1)} \gets \theta_\mathrm{sh}^{(t)} + \eta_g (\tilde\Delta_\mathrm{sh}^{(t+1)} + \xi)$, $\xi \sim \mathcal{N}(0, \sigma^2 I)$
\State $\theta_\mathrm{cl}^{(t+1)} \gets \theta_\mathrm{cl}^{(t)} + \eta_g \overline{\Delta_\mathrm{cl}}$ \Comment{head: FedAvg, no noise}
\end{algorithmic}
\end{algorithm}

\begin{figure}[t]
\centering
\resizebox{\linewidth}{!}{%
\begin{tikzpicture}[
  font=\scriptsize,
  >={Latex[length=1.6mm]},
  thick,
  box/.style        = {draw, rounded corners=2pt, align=center,
                        inner sep=2.5pt, font=\scriptsize},
  client/.style     = {box, fill=green!15, minimum width=11mm, minimum height=6mm},
  server/.style     = {box, fill=blue!12, minimum width=16mm, minimum height=10mm},
  feat/.style       = {box, fill=blue!25, minimum width=14mm, minimum height=5mm},
  head/.style       = {box, fill=orange!30, minimum width=14mm, minimum height=5mm},
  summary/.style    = {box, fill=yellow!40, minimum width=11mm, minimum height=5mm},
  authority/.style  = {box, fill=red!15, minimum width=17mm, minimum height=9mm},
  certbox/.style    = {box, fill=red!8, minimum width=35mm, minimum height=5mm},
  update/.style     = {->, thick},
  dupdate/.style    = {->, thick, dashed, gray},
  zone/.style       = {draw=gray!40, dashed, rounded corners=3pt, inner sep=3mm},
]

% -------- Top zone: TRAINING TIME --------
\node[client] (cA) at (0, 0)       {Client $i$};
\node[server] (srv1) at (3.3, 0)   {Server \\ \scriptsize (aggregator)};
\node[feat]   (fb)   at (6.6, 0.5) {Features $\theta_\mathrm{sh}$};
\node[head]   (hd)   at (6.6,-0.5) {Head $\theta_\mathrm{cl}$};

\draw[update] (cA.east) -- node[above,font=\tiny] {$\Delta_i$} (srv1.west);
\draw[update] (srv1.east |- fb.west) -- node[above,font=\tiny] {EMA $\alpha$} (fb.west);
\draw[update] (srv1.east |- hd.west) -- node[below,font=\tiny] {FedAvg} (hd.west);

\node[summary] (sum) at (1.7, -1.5) {$\sigma_i\!+\!\eta$};
\draw[update] (cA.south) |- (sum.west);
\node[font=\tiny,gray!60!black] at (2.9, -1.85) {hash + commit};
\draw[update] (sum.east) -| ($(srv1.south) + (-2mm,0)$);

\node[anchor=west,font=\scriptsize,gray!70!black] at (-0.7, 1.0)
  {\textbf{Training time}};

\begin{pgfonlayer}{background}
  \node[zone, fit=(cA)(sum)(fb)(hd)] {};
\end{pgfonlayer}

% -------- Bottom zone: REQUEST TIME --------
\node[client]    (cB)   at (0, -3.6)   {Client $i^\star$};
\node[font=\tiny] at (0, -4.05) {\textit{``forget me''}};
\node[authority] (auth) at (3.3, -3.6) {Authority \\ \scriptsize + Server};
\node[summary]   (op)   at (6.0, -3.6) {open $\sigma_{i^\star}$};
\node[head]      (hd2)  at (8.5, -3.6) {$\theta_\mathrm{cl}'$};

\draw[update] (cB.east) -- node[above,font=\tiny] {req} (auth.west);
\draw[update] (auth.east) -- (op.west);
\draw[update] (op.east)  -- node[above,font=\tiny]
  {$\gamma\,\Pi_{B(0,R)}(v)$} (hd2.west);

\node[certbox] (cert) at (5.0, -4.9)
  {certificate $(\mu_\mathrm{joint}(T),\,\delta,\,h_\sigma,\,\mathrm{tr})$};
\draw[update] (hd2.south) |- (cert.east);
\draw[dupdate] (cert.west) -| node[below,font=\tiny,gray] {verify} (auth.south);

\node[anchor=west,font=\scriptsize,gray!70!black] at (-0.7, -2.65)
  {\textbf{Request time}};

\begin{pgfonlayer}{background}
  \node[zone, fit=(cB)(auth)(hd2)(cert)] {};
\end{pgfonlayer}

\end{tikzpicture}%
}
\caption{FORGE at a glance. \emph{Top (training time):} each client sends a
clipped update $\Delta_i$ into the decayed aggregator and, every $F$ rounds, a
committed DP summary $\sigma_i$. The feature block is updated by an EMA with
decay $\alpha$ and server noise; the head is updated by unweighted FedAvg.
\emph{Bottom (request time):} the departing client's commitment is opened, the
server applies a bounded head-only step, and an audit certificate carrying the
composed GDP parameter and the hash-chain transcript is emitted for the
authority to verify.}
\label{fig:forge-schematic}
\end{figure}

\begin{algorithm}[t]
\footnotesize
\caption{FORGE non-interactive erasure.}
\label{alg:erasure}
\begin{algorithmic}[1]
\Require $(i^\star, r_\forget, r_\mathrm{req})$; $\theta^{(r_\mathrm{req})}$; $\sigma_{i^\star}$; $\gamma$, $R$
\State $T \gets r_\mathrm{req} - r_\forget$; verify $\mathrm{hash}(\sigma_{i^\star})$
\State $v \gets \nabla_{\theta_\mathrm{cl}} \mathcal{L}_\mathrm{sub}(\theta_\mathrm{cl}^{(r_\mathrm{req})}; \sigma_{i^\star})$
\State $\theta_\mathrm{cl}^{(r_\mathrm{req})} \gets \theta_\mathrm{cl}^{(r_\mathrm{req})} + \gamma \Pi_{B(0, R)}(v)$
\State $\mathrm{cert} \gets (\mu_\mathrm{joint}(T), \delta_\summary, \mathrm{hash}(\sigma_{i^\star}))$ \Comment{$\mu_\mathrm{joint}(T) {=} \sqrt{\mu_1(T)^2 {+} \mu_\summary^2}$}
\State Purge $\sigma_{i^\star}$; \Return $\theta^{(r_\mathrm{req})}, \mathrm{cert}$
\end{algorithmic}
\end{algorithm}

\subsection{Design remarks}
\label{sec:method-remarks}

We briefly comment on three design decisions. First, the head is the smallest module
whose output distribution can be pushed off a client's memorised instances by a bounded
parameter change. Operating on the head rather than the full model confines the
perturbation to a small subspace, limits utility damage, and establishes the
independence-of-noise condition that Theorem~\ref{thm:composition} relies on; if the
subtraction were applied to the whole model, the summary mechanism would overwrite the
feature-block release and the adaptive GDP composition would collapse to sequential
composition on correlated releases. Second, if summaries were used
during training to influence the model, they would become part of the training
trajectory and the forward-forgetting rate's sensitivity analysis would need to account
for their contribution. The summary channel is therefore designed to treat summaries as
inert during training; they are accessed only at erasure. This keeps the training-time
analysis (Theorem~\ref{thm:forward-forgetting}) separate from the erasure-time analysis
(Lemma~\ref{lem:subtraction}) and in turn permits adaptive Gaussian-DP composition. Third, the
feature decoupling in Algorithm~\ref{alg:decayed-aggregation} differs from standard
FedAvg in two ways: (i) the EMA on $\tilde\Delta_\mathrm{sh}$ slows the feature-block's
effective learning rate by factor $(1-\alpha)$, and (ii) the server noise $\xi$ perturbs
the feature block each round. Section~\ref{sec:results-utility} measures the empirical
cost of these modifications relative to the vanilla-FedAvg retrain baseline on
CIFAR-10.
